\chapter{Conclusione}
\label{ch:conclusione}
Il lavoro svolto ha consentito di sostituire una gestione caratterizzata da una digitalizzazione parziale del processo, nella quale il sistema preesistente era limitato all'archiviazione dei campioni analizzati e dei risultati prodotti, mentre le fasi operative venivano svolte e documentate manualmente. ENVASS 2.0 introduce una piattaforma integrata in grado di supportare e tracciare l'intero flusso di lavoro del laboratorio, garantendo una gestione più strutturata delle attività e una maggiore affidabilità, coerenza e verificabilità delle informazioni gestite.

Gli obiettivi del progetto sono stati raggiunti attraverso l’introduzione di soluzioni architetturali volte a migliorare la gestione dei dati e dei processi del sistema. In particolare:

\begin{itemize}
	\item \textbf{La centralizzazione delle informazioni} ha consentito una riduzione delle ridondanze e una maggiore uniformità nella gestione dei dati anagrafici e operativi;
	\item \textbf{L'automazione del workflow} ha reso più coerente la transizione tra le diverse fasi del processo analitico e la generazione della documentazione associata;
	\item \textbf{I meccanismi di immutabilità dei dati validati} hanno garantito il rispetto dei requisiti di integrità e tracciabilità, impedendo modifiche non controllate;
	\item \textbf{Il sistema di audit trail} ha permesso la registrazione strutturata delle operazioni effettuate, migliorando la verificabilità delle attività.
\end{itemize}

Nel complesso, la piattaforma ENVASS 2.0 rappresenta un'evoluzione significativa rispetto alla precedente versione, introducendo un modello informativo più coerente, tracciabile e aderente ai requisiti di qualità richiesti dal contesto di laboratorio.
Poiché il sistema è stato progettato per essere adottato operativamente nel breve periodo, la valutazione del suo impatto su larga scala richiederà un utilizzo prolungato in condizioni reali di esercizio.

Con il rilascio della presente release si conclude la fase di sviluppo attivo del progetto. L'evoluzione della piattaforma e l'ottimizzazione progressiva delle funzionalità vengono pertanto demandate alle attività di manutenzione ordinaria e al supporto operativo, guidati dai feedback che emergeranno dagli utenti finali durante l'effettivo utilizzo sul campo.

In ottica di sviluppi futuri, il sistema potrà essere ulteriormente esteso attraverso l’integrazione di funzionalità basate sull’intelligenza artificiale, finalizzate al supporto delle attività analitiche e decisionali. In particolare, potranno essere introdotti moduli per l’analisi automatica dei dati visivi, ad esempio per il conteggio delle colonie nelle piastre Petri, con l’obiettivo di ridurre il carico di lavoro manuale e migliorare la riproducibilità dei risultati. Ulteriori estensioni potrebbero riguardare sistemi di supporto alla modifica e all’elaborazione intelligente delle immagini, nonché l’introduzione di notifiche automatiche alla conclusione del periodo di incubazione, migliorando così l’efficienza complessiva dei flussi operativi del laboratorio